\subsection{Invariant Factors}

\begin{definition} \label{def:operator-invariant-factors-elementary-divisors} \hlt{Invariant Factors and Elementary Divisors}\\
Let $0<\dim\mathcal{V}<\infty$, $T\in\LL(\mathcal{V})$. Then $\mathcal{V}_T$ is a finitely generated torsion $\mathbb{F}[z]$-module by Proposition \ref{prop:operator-modules-are-fg-torsion}(i)--(ii).
\begin{enumerate}[label=(\roman*)]
\setlength{\itemsep}{0pt}
\item For $r\in\N$, a \hlt{cyclic decomposition} of $T$ is $\mathcal{V}=\bigoplus\limits_{j=1}^r\mathcal{C}_T(\mathbf{v}_j)$ with $(\mathbf{v}_1,\ldots,\mathbf{v}_r)\in(\mathcal{V}\setminus\{\mathbf{0}\})^r$.
\item The \hlt{invariant factors} of $T$ are the unique monic representatives of $\operatorname{Inv}(\mathcal{V}_T)$ from Definition \ref{def:invariant-factors-pid-module}.\\
Write $\operatorname{InvFact}(T)=(f_1,\ldots,f_r)$. Then each $f_j$ is monic and nonconstant, and $f_r\mid\cdots\mid f_1$.
\item The \hlt{elementary divisors} of $T$ are the unique monic representatives of $\operatorname{ElemDiv}(\mathcal{V}_T)$ from Definition \ref{def:elementary-divisors-pid-module}.\\
Denote the resulting multiset by $\operatorname{ElemDiv}(T)$.
\item Let $\mathcal{I}$ be finite, $(\mathbf{v}_{\alpha})_{\alpha\in\mathcal{I}}\in(\mathcal{V}\setminus\{\mathbf{0}\})^{\mathcal{I}}$, $\mathcal{V}=\bigoplus\limits_{\alpha\in\mathcal{I}}\mathcal{C}_T(\mathbf{v}_{\alpha})$. For $\alpha\in\mathcal{I}$, set $d_{\alpha}\defeq\deg m_{\mathbf{v}_{\alpha}}$ and $\mathcal{B}_{\alpha}\defeq(\mathbf{v}_{\alpha},T\mathbf{v}_{\alpha},\ldots,T^{d_{\alpha}-1}\mathbf{v}_{\alpha})$.\\
After choosing an ordering of $\mathcal{I}$, the corresponding concatenation of the $\mathcal{B}_{\alpha}$ is an \hlt{elementary divisor basis}\\
if $\{m_{\mathbf{v}_{\alpha}}:\alpha\in\mathcal{I}\}=\operatorname{ElemDiv}(T)$ as multisets.
\item For $n\geq1$ and $A\in\operatorname{Mat}(n,n,\mathbb{F})$, set $\operatorname{InvFact}(A)\defeq\operatorname{InvFact}(T_A)$. Set $\operatorname{ElemDiv}(A)\defeq\operatorname{ElemDiv}(T_A)$.
\end{enumerate}
\end{definition}

\begin{theorem} \label{thm:operator-invariant-factor-skeleton} \hlt{Invariant-Factor Skeleton for Operators}\\
Let $0<\dim\mathcal{V}<\infty$, $T\in\LL(\mathcal{V})$, $r\in\N$, and $\operatorname{InvFact}(T)=(f_1,\ldots,f_r)$.
\begin{enumerate}[label=(\roman*)]
\setlength{\itemsep}{0pt}
\item $(\exists\,\mathbf{v}_1,\ldots,\mathbf{v}_r\in\mathcal{V}\setminus\{\mathbf{0}\})$ such that, for $\mathcal{C}_j\defeq\mathcal{C}_T(\mathbf{v}_j)$, $\mathcal{V}=\bigoplus\limits_{j=1}^r\mathcal{C}_j$. Then $m_{\mathbf{v}_j}=f_j$, $\mathcal{C}_j\cong\mathbb{F}[z]/\langle f_j\rangle$ for $j\in\{1,\ldots,r\}$.
\item Set $\mathcal{B}_j\defeq(\mathbf{v}_j,T\mathbf{v}_j,\ldots,T^{\deg f_j-1}\mathbf{v}_j)$ and $\mathcal{B}\defeq(\mathcal{B}_1,\ldots,\mathcal{B}_r)$. Then $\mathcal{B}$ is a basis of $\mathcal{V}$ and $[T]_{\mathcal{B}\leftarrow\mathcal{B}}=\bigoplus\limits_{j=1}^rC_{f_j}$.
\item $\mu_T=f_1$.
\item $q_T=\prod\limits_{j=1}^rf_j$.
\end{enumerate}
\begin{proof}
\begin{enumerate}[label=(\roman*)]
\setlength{\itemsep}{0pt}
\item The module $\mathcal{V}_T$ is finitely generated and torsion by Proposition \ref{prop:operator-modules-are-fg-torsion}(i)--(ii).\\
Choose an $\mathbb{F}[z]$-module isomorphism $\Psi:\bigoplus\limits_{j=1}^r\mathbb{F}[z]/\langle f_j\rangle\to\mathcal{V}_T$ by Theorem \ref{thm:pid-invariant-factor-decomposition}(i).\\
Let $\mathbf{e}_j$ have $j$th coordinate $1+\langle f_j\rangle$ and all other coordinates zero. Set $\mathbf{v}_j\defeq\Psi(\mathbf{e}_j)$ and $\mathcal{C}_j\defeq\Psi(\mathbb{F}[z]\mathbf{e}_j)$.\\
Since $f_j$ is nonconstant and $\Psi$ is injective, $\mathbf{e}_j\neq\mathbf{0}$ and $\mathbf{v}_j\neq\mathbf{0}$. For $p\in\mathbb{F}[z]$, $\Psi(p\mathbf{e}_j)=p(T)\mathbf{v}_j$ by $\mathbb{F}[z]$-linearity.\\
$\mathcal{C}_j=\mathcal{C}_T(\mathbf{v}_j)$ by the definitions of $\mathcal{C}_j$ and $\mathcal{C}_T(\mathbf{v}_j)$. Since $\Psi$ is bijective, $\mathcal{V}=\bigoplus\limits_{j=1}^r\mathcal{C}_j$ and $\mathcal{C}_j\cong\mathbb{F}[z]/\langle f_j\rangle$.\\
The annihilator of $\mathbf{e}_j$ is $\langle f_j\rangle$, so $\operatorname{ann}_T(\mathbf{v}_j)=\langle f_j\rangle$ by $\Psi$. Hence $m_{\mathbf{v}_j}=f_j$ by Definition \ref{def:vector-order-under-operator}.
\item For $j\in\{1,\ldots,r\}$, $\mathcal{B}_j$ is a basis of $\mathcal{C}_j$ by Theorem \ref{thm:cyclic-subspace-quotient-and-companion}(ii). $[T|_{\mathcal{C}_j}]_{\mathcal{B}_j\leftarrow\mathcal{B}_j}=C_{f_j}$ by Theorem \ref{thm:cyclic-subspace-quotient-and-companion}(iii).\\
The concatenation $\mathcal{B}$ is a basis of $\mathcal{V}$ by part (i) and Definition \ref{def:adapted-basis-direct-sum}.\\
Since $(\forall\,j\in\{1,\ldots,r\})\,T\mathcal{C}_j\subseteq\mathcal{C}_j$, all off-diagonal blocks vanish. Thus $[T]_{\mathcal{B}\leftarrow\mathcal{B}}=\bigoplus\limits_{j=1}^rC_{f_j}$.
\item For $j\in\{1,\ldots,r\}$, $\operatorname{Ann}_{\mathbb{F}[z]}(\mathcal{C}_j)=\langle f_j\rangle$ because $\mathcal{C}_j=\mathbb{F}[z]\mathbf{v}_j$ and $m_{\mathbf{v}_j}=f_j$ by part (i).\\
$\operatorname{Ann}_{\mathbb{F}[z]}(\mathcal{V}_T)=\langle\operatorname{lcm}(f_1,\ldots,f_r)\rangle$ by part (i), Proposition \ref{prop:orders-direct-sum-lcm}. Since $f_r\mid\cdots\mid f_1$, $\operatorname{lcm}(f_1,\ldots,f_r)=f_1$ and $\operatorname{Ann}_{\mathbb{F}[z]}(\mathcal{V}_T)=\langle f_1\rangle$.\\
$\operatorname{Ann}_{\mathbb{F}[z]}(\mathcal{V}_T)=\langle\mu_T\rangle$ by Proposition \ref{prop:minimal-polynomial-module-order}(i). Since $f_1$ and $\mu_T$ are monic, $f_1=\mu_T$.
\item Set $A\defeq[T]_{\mathcal{B}\leftarrow\mathcal{B}}$. Then $q_T=p_A$ by Definition \ref{def:operator-characteristic-polynomial}.\\
$A=\bigoplus\limits_{j=1}^rC_{f_j}$ by part (ii), so $p_A=\prod\limits_{j=1}^rp_{C_{f_j}}$ by repeated application of Proposition \ref{prop:block-upper-triangular-characteristic-polynomial}(i).\\
For $j\in\{1,\ldots,r\}$, $p_{C_{f_j}}=f_j$ by Theorem \ref{thm:companion-matrix-has-prescribed-minimal-polynomial}(ii). Hence $q_T=\prod\limits_{j=1}^rf_j$.
\end{enumerate}
\hlt{Structural Commentary}: PID quotient generators produce cyclic bases; annihilators and determinants recover $\mu_T$ and $q_T$.\\
\hlt{Key Insight}: Cyclic quotient generators turn invariant factors into companion blocks.
\end{proof}
\end{theorem}
